\section{Conclusion}
\label{sec:conclusion}

We presented a PPO-scheduled, LoRA-adapted latent diffusion pipeline for visual counterfactual explanation on chest radiographs, with a soft background lock that largely confines edits to a dilated lung region. We evaluated it under a protocol that pairs the field's conventional classifier-scored validity with a lesion-annotation-anchored hit rate and a held-out transfer ratio. The system attains a 2.60$\times$ lesion-targeting lift over chance, and its ensemble-guided policy improves both lesion targeting and mean held-out transfer over a single-classifier policy whose configuration differs only in the guidance classifier, in a comparison of one training run per arm on partly different evaluation images. It also shows, on its own outputs, that high CV coexists with diffuse edits and adversarial-regime transfer on the most independent validator. Two mechanisms designed to concentrate edits were trained and measured rather than assumed, and both measured as harmful in a single training run each, one by the largest effect size in the study. We take the broader implication to be methodological: CV computed with the guidance classifier should not be reported as a standalone success criterion, and an anchored companion metric (to human annotation, to model independence, or both) costs little and changes the conclusion.
